\section*{Chapter \ref{ch:iv1:timed-test-10} --- Timed Test 10}
\markboth{\protect\omTocMark{}{Answers: Timed Test 10}}{\protect\omTocMark{}{Answers: Timed Test 10}}

\begin{solution}{iq:iv1:timed-test-10:1}
The probability is $\frac23$.

The sum falls below 1 inside the corner simplex of the cube, which has volume $\frac{1}{3!} = \frac16$\index{volume of a simplex}. By the symmetry $x \to 1 - x$ for all three, the sum exceeds 2 with the same chance. So the middle has $1 - \frac26 = \frac23$.

A simulation returns $0.667$. Here, the mean is $1.5$ and the standard deviation is $\frac12$, so the interval runs one standard deviation either side, and $\frac23$ is close to the normal rule of thumb of $0.68$.

\iqwrong{Treating the sum as uniform on $[0, 3]$, which gives $\frac13$. Sums bunch towards the middle.}
\end{solution}

\begin{solution}{iq:iv1:timed-test-10:2}
The probability is $\frac{50}{81} \approx 0.617$.

There are $3^5 = 243$ outcomes. Use inclusion--exclusion\index{inclusion--exclusion} on the empty bins. With a given bin empty, there are $2^5 = 32$ outcomes, and with two given bins empty, there is 1. So the outcomes with no empty bin are
\[
243 - 3 \times 32 + 3 \times 1 = 150.
\]
That gives $\frac{150}{243} = \frac{50}{81}$.

A second count uses the sizes: $(3, 1, 1)$ has $3 \times \binom53 \times 2 = 60$ outcomes, and $(2, 2, 1)$ has $3 \times \binom52\binom32 = 90$. The total is 150. An enumeration agrees.

\iqwrong{Forgetting the last term and getting $\frac{147}{243}$. The outcomes with two empty bins are subtracted twice each in the single-bin count, so they must be added back.}
\end{solution}

\begin{solution}{iq:iv1:timed-test-10:3}
Chebyshev gives at most $0.16$, and the normal value is about $0.0124$.

Chebyshev says that $P(|X - \mu| \ge a) \le \frac{\sigma^2}{a^2}$\index{Chebyshev's inequality}, which is $\frac{4}{25} = 0.16$.

For a normal, the deviation is $\frac52 = 2.5$ standard deviations. Both tails count: $2 \times 0.0062 = 0.0124$. A simulation agrees.

The gap, of a factor of 13, is the cost of assuming nothing about the shape. And Chebyshev's bound is attained by a three-point distribution, which puts all the deviation exactly at $\pm5$. So it cannot be improved without an assumption.

Both are correct: only the first holds for any distribution with this variance, and this matters when a tail is fat.

\iqwrong{Using only one tail for the normal, $0.0062$. The question asks for a deviation in either direction.}
\end{solution}

\begin{solution}{iq:iv1:timed-test-10:4}
Each probability is $\frac59$.

Each die has three values, each with chance $\frac13$. So there are nine equally likely pairs.

A against B: 2 beats 1, 4 beats 1, and 9 beats all three. That is $1 + 1 + 3 = 5$ pairs out of 9.

B against C: 6 beats 3 and 5, and 8 beats all three. That is $2 + 3 = 5$.

C against A: 3 beats 2, 5 beats 2 and 4, and 7 beats 2 and 4. That is $1 + 2 + 2 = 5$.

So A beats B, B beats C and C beats A. An enumeration gives the same. The relation ``beats'' is not transitive\index{non-transitive relation}, although each die has the same mean of 5. So a choice made second always has the edge.

\iqwrong{Assuming that one die must be best, because a ranking looks natural. Ranking by one-on-one chance need not give an order.}
\end{solution}

\begin{solution}{iq:iv1:timed-test-10:5}
The probability is $\frac23$, and the expected wait is $\frac{20}{3} \approx 6.67$ minutes.

The rates are $\frac{1}{10}$ and $\frac{1}{20}$. By the memoryless property, the wait for each is exponential, whatever has happened before you arrive. The first bus is on line A with chance\index{exponential race} $\frac{1/10}{1/10 + 1/20} = \frac23$.

The first arrival has rate $\frac{3}{20}$, so its mean is $\frac{20}{3}$. A simulation agrees.

Limit check: if line B is removed, its rate is 0, the chance is 1, and the wait is 10. So the wait is below the shorter gap, which makes sense.

The choice between lines and the waiting time are independent. So the wait is $\frac{20}{3}$ whichever line you end up on.

\iqwrong{Taking half of each gap as the wait, which is right only for regular buses. A Poisson line has a mean wait equal to its whole mean gap.}
\end{solution}

\begin{solution}{iq:iv1:timed-test-10:6}
The smallest group has 6 people.

The chance that all $n$ birthdays differ\index{birthday problem} is $\prod_{i = 0}^{n - 1}(1 - \frac{i}{20})$. At $n = 5$, this is $\frac{20 \cdot 19 \cdot 18 \cdot 17 \cdot 16}{20^5} = 0.581$. At $n = 6$, it is $0.581 \times \frac{15}{20} = 0.436$. So the chance of a match goes from $0.419$ to $0.564$, and it crosses $\frac12$ at 6.

The rule of thumb is $1.18\sqrt{d}$: with $d = 20$, it is $5.3$, which rounds up to 6. An exact loop confirms it. The same rule gives $22.5$ for 365 days, near the true 23.

\iqwrong{Answering 10, half the days, because half seems natural. Pairs grow as $n^2$, so the answer scales with the square root.}
\end{solution}

\begin{solution}{iq:iv1:timed-test-10:7}
The expected value is $\frac{0.5}{\ln 2} \approx 0.721$.

The joint density is $\frac1x$ for $0 < y < x < 1$. So given $Y = y$, the density of $X$ is proportional to $\frac1x$ on $(y, 1)$\index{conditional density}. The constant is $-\ln y$. Then
\[
\mathbb{E}[X \mid Y = y] = \frac{\int_y^11\,dx}{-\ln y} = \frac{1 - y}{-\ln y}.
\]
At $y = \frac12$, this is $\frac{0.5}{0.693} = 0.721$.

Check: it lies between $\frac12$ and 1, as it must. As $y \to 1$, the formula tends to 1. As $y \to 0$, it tends to 0, which is correct, as a small $Y$ points to a small $X$. A simulation in a thin window around $\frac12$ agrees.

\iqwrong{Answering $\frac34$, the midpoint of the allowed range. The factor $\frac1x$ favours small $X$, because a small $X$ puts more density at any given $Y$.}
\end{solution}

\begin{solution}{iq:iv1:timed-test-10:8}
With five dice the probability is exactly $\frac12$, and with six dice it is about $0.546$.

Replace each face $x$ by $7 - x$, which leaves every die's law unchanged\index{symmetry}. With five dice, the total $T$ becomes $35 - T$, so $T$ and $35 - T$ have the same law. Now $T \ge 18$ is the same event as $35 - T \le 17$. So $P(T \ge 18) = P(T \le 17)$, and these two events cover everything, with no overlap. Each therefore has chance $\frac12$.

With six dice, $T$ and $42 - T$ have the same law, and the centre, 21, is a possible total. So $P(T > 21) = P(T < 21)$, and
\[
P(T \ge 21) = \frac12 + \frac12P(T = 21) = \frac12 + \frac12 \cdot \frac{4\,332}{46\,656} \approx 0.546.
\]

\iqwrong{Giving $\frac12$ for six dice as well. Five dice have a centre of $17.5$, which no total can equal, but six dice have a centre they can hit, and ``at least'' counts it.}
\end{solution}

\begin{solution}{iq:iv1:timed-test-10:9}
The third item is kept with probability $\frac15$. For a sample of 2, keep each new item with chance $\frac2i$, replacing one of the two at random, and each item then has chance $\frac25$.

The third item enters with chance $\frac13$. It survives the fourth with $\frac34$ and the fifth with $\frac45$\index{reservoir sampling}. The product is $\frac13 \cdot \frac34 \cdot \frac45 = \frac15$.

The product telescopes for any $i$ and $n$, so every item ends up with chance $\frac1n$. You never need to know $n$ in advance.

For a sample of $k$, item $i > k$ enters with chance $\frac{k}{i}$. Each later item $j$ removes it with chance $\frac{k}{j} \cdot \frac1k = \frac1j$, so it survives with $\frac{j - 1}{j}$. The product is $\frac{k}{n}$. A simulation confirms the uniform shares.

\iqwrong{Using a fixed chance per item, say $\frac{1}{100}$. It needs the length, and it gives later items a better share.}
\end{solution}

\enlargethispage{-6\baselineskip}%
\begin{solution}{iq:iv1:timed-test-10:10}
The probability is about $0.0211$.

Choose the rank: 13. Choose its three suits\index{counting hands}: $\binom43 = 4$. Choose the two other ranks: $\binom{12}{2} = 66$. Choose a suit for each: $4^2 = 16$. The count is
\[
13 \times 4 \times 66 \times 16 = 54\,912,
\]
out of the total, so $0.0211$.

An independent simulation returns the same value. Note that the two ranks must differ, and if they did not, the hand would be a full house.

\iqwrong{Using 49 choices for the fourth card and 48 for the fifth. This counts full houses, and it gives each hand twice.}

\iqfollow{What is the chance of a full house?}{$13 \times 4 \times 12 \times 6 = 3744$ hands, or $0.00144$. It is rarer, even though the description is similar.}
\end{solution}

\begin{solution}{iq:iv1:timed-test-10:11}
The probability is about $0.057$.

The count has mean 50 and standard deviation 5\index{continuity correction}. We want $H \le 40$ or $H \ge 60$. The correction puts the edges at $40.5$ and $59.5$, each $1.9$ standard deviations from the mean. Two tails give $2 \times 0.0287 = 0.0574$.

An exact binomial sum gives $0.0569$, so the correction is very close. Without the correction, $z = 2$, and the answer would be $0.0455$. This is off by a fifth, a large error for such a simple step.

A check against a known figure: two standard deviations hold about $95\%$, so about $5\%$ lie beyond. A count that can fall exactly on the edge adds a little more.

\iqwrong{Using one tail only, which gives $0.029$, half of the answer. The question asks for a deviation either way.}
\end{solution}

\begin{solution}{iq:iv1:timed-test-10:12}
The put should cost 15. At 14, buy the put and the stock, and sell the call, for a certain profit of 1.

Put--call parity says that $C - P = S - K$ at a zero rate\index{put--call parity}. So $P = C - S + K = 10 - 100 + 105 = 15$.

At 14, the put is too cheap. Buy the put, buy the stock and sell the call. The cost is $14 + 100 - 10 = 104$. At expiry, the stock plus the put minus the call is worth exactly $K = 105$ in every case: if $S_T > 105$, the call is exercised against you at 105, and if $S_T < 105$, you exercise the put at 105. So the profit is 1, with no risk.

A check: at $S_T = 90$ and at $S_T = 120$, the position is worth 105.

\iqwrong{Selling the put, because it looks like a poor asset at 14. The rule is to buy what is cheap against its replication.}
\end{solution}
